% !TeX root = ../Book2.tex
\chapter*{序言}
\addcontentsline{toc}{chapter}{序言}
把一个空间切成简单的部分，能否从这些部分及其粘合方式恢复整体信息？在空间上选取映射、向量场或向量丛时，又有哪些选择能够连续变形，哪些选择受到整体拓扑的限制？本卷围绕这两类问题讨论代数拓扑与微分拓扑。空间的分解给出可计算的代数对象，光滑映射的逆像、交点和临界点则让这些代数对象重新获得几何解释。

\ProseParagraphBreak

圆周绕行的次数，是这种联系的一个早期例子。对于圆周上以固定点为起讫点的闭路，一个整数可以记录其同伦类；把闭路看成圆周的自映射，这个整数也是它在一维同调上诱导的乘数。在适当的光滑与定向条件下，它还可以通过正则值逆像点的带符号计数得到。本卷会多次从不同方法回到同一个对象。比较这些构造时，需要说明它们怎样对应、依赖哪些选择，以及为什么得到相同的不变量。

\ProseParagraphBreak

本卷面向已有基本拓扑、线性代数和群论训练的数学系本科生与研究生。开始代数拓扑部分，需要理解连续映射、开覆盖、积与商拓扑、紧性，以及群、子群和商群。\chapref{g2:ch:01}重新建立本卷所需的基本群与覆盖理论，并引入胞腔空间，读者可以由此直接开始。链复形、正合列、系数变换等代数工具随构造展开；熟悉核、像、商和线性映射，有助于看清这些工具正在记录什么。

\ProseParagraphBreak

第一编由空间分解建立同调，再加入系数、乘积和对偶。这里既要把定义做成可用的计算，也要解释不同模型为何给出相容的结果。局部系统记录沿环路移动后发生的系数变化，流形对偶联系不同维数的同调信息，微分形式与层上同调的比较则把连续拓扑、积分和局部粘合联系起来。第二编进一步研究高阶同伦和纤维化，利用逐层逼近组织计算；谱序列、障碍与上同调运算各自解决不同的问题，在所需构造建立之后依次出现。

\ProseParagraphBreak

第三编研究光滑映射与流形。横截性使逆像具有可控制的局部结构，相交与映射度将局部符号汇成整体计数，临界点和把手则给出流形的分解。第四编转向向量丛：先研究同构分类，再构造示性类，最后用联络的曲率写出相应的微分形式。第五编讨论配边、稳定化与 K 理论，把流形的几何关系和向量丛的运算放到新的分类问题中，并引向椭圆算子的指标。

\ProseParagraphBreak

这些内容可以按问题分别学习。初学代数拓扑时，可从\chapref{g2:ch:01}到\chapref{g2:ch:04}建立同伦、同调与乘积的基本经验，再学习局部系数、对偶及同伦计算。希望先研究光滑映射者，可在掌握第一卷第11、12章的流形、局部标准形和单位分解后，从本卷\chapref{g2:ch:14}开始；邻域与局部平凡性还要用到第一卷第13章的流。随后把几何计数同同调配对比较、把临界点分解同胞腔同调联系时，需要回到本卷第一编的相关内容。微分形式、积分和联络分别可查第一卷第14、15章。

\ProseParagraphBreak

以向量丛为主要兴趣的读者，可结合\chapref{g2:ch:01}的胞腔构造阅读\chapref{g2:ch:19}，并按需回顾第一卷的局部平凡化与单位分解。继续学习示性类，需要第一编的上同调乘积和相对理论；遇到首障碍和平方运算时，还需查阅\chapref{g2:ch:12}与\chapref{g2:ch:13}的相关部分。学习曲率代表，则需第一卷第14、18章的微分形式和主丛联络。各编导读和各章开头进一步说明具体所需结果。星号标记的章节与习题、以及附录中的专题，供读者按目的选读。

\ProseParagraphBreak

阅读本卷时，宜让抽象构造始终伴随一个具体计算。先在球面、环面或射影空间上写出链群和边界映射，再理解自然性与比较定理；先追踪一张纤维化中基空间、纤维和总空间的信息，再逐页阅读谱序列。证明中的辅助选择同样值得留意：定向、基点、局部平凡化、横截扰动都会进入中间表达，所得对象在哪一种等价意义下不变，需要由相应论证来回答。章末习题及解答用于训练和检查这些判断。

\ProseParagraphBreak

附录分别延伸结与低维拓扑、分层空间、等变与稳定方法、叶状结构与映射奇点，以及有理同伦。正文与附录的阅读说明给出可继续学习的教材和专题文献，具体引用标明外部结果的版本与用途。书中给出证明的构造可以沿正文逐步核验；以外部定理使用的结果，则在相应位置说明条件和来源。例如第五编具体展开复 K 理论的周期性证明，实 KO 理论的八周期与一般指标定理采用注明来源的结果。不同资料承担的任务由其使用位置说明。

\ProseParagraphBreak

本卷所得的上同调、相交、示性类与稳定方法，将用于后续各卷的几何与分析问题。读者可以在完成相关部分后继续学习，也可以在遇到具体问题时回查相应构造。当前二十四章、附录及书前书后文字均为初稿，尚未经过正式审核。
